\documentclass[a4paper,11pt]{article}
\usepackage[utf8]{inputenc}
\usepackage[T1]{fontenc}
\usepackage[german]{babel}
\usepackage{amsmath,amssymb,amsfonts}
\usepackage{bm}
\usepackage{geometry}
\usepackage{microtype}
\usepackage{booktabs}
\usepackage{tabularx}
\usepackage{array}
\usepackage{listings}
\usepackage{xcolor}
\usepackage{xurl}
\usepackage[round,authoryear]{natbib}
\usepackage{hyperref}

\geometry{margin=2.5cm}

\hypersetup{
    colorlinks=true,
    linkcolor=blue,
    citecolor=blue,
    urlcolor=blue
}

% Mathe-Makros in PDF-Lesezeichen (hyperref) neutralisieren
\pdfstringdefDisableCommands{%
  \def\bm#1{#1}%
  \def\tfrac#1#2{#1/#2}%
  \def\tilde#1{#1}%
}

\definecolor{codegreen}{rgb}{0,0.6,0}
\definecolor{codegray}{rgb}{0.5,0.5,0.5}
\definecolor{codepurple}{rgb}{0.58,0,0.82}
\definecolor{backcolour}{rgb}{0.96,0.96,0.94}

\lstdefinestyle{mystyle}{
    backgroundcolor=\color{backcolour},
    commentstyle=\color{codegreen},
    keywordstyle=\color{magenta},
    numberstyle=\tiny\color{codegray},
    stringstyle=\color{codepurple},
    basicstyle=\ttfamily\footnotesize,
    breakatwhitespace=false,
    breaklines=true,
    captionpos=b,
    keepspaces=true,
    numbers=left,
    numbersep=5pt,
    showspaces=false,
    showstringspaces=false,
    showtabs=false,
    tabsize=2
}
\lstset{style=mystyle}

\newcommand{\Ntilde}{\tilde{N}}
\newcommand{\code}[1]{\texttt{\small #1}}
\newcommand{\file}[1]{\texttt{\small #1}}

% Saubere Ränder: keine microtype-Protrusion (Bindestriche/Klammern/Schrägstriche
% ragen sonst optisch in den Rand); Notfall-Dehnung gegen überstehende Zeilen.
\microtypesetup{protrusion=false}
\emergencystretch=3em
\hyphenation{Con-straint}

% Große Abschnitte (\section) immer auf einer neuen Seite beginnen.
\let\stdsection\section
\renewcommand{\section}{\clearpage\stdsection}

\title{\textbf{Reibungsfreier Mortar-Kontakt in EdelweissFE}\\
\large Theorie, Herleitung und Implementierung des Normalkontakts\\
\normalsize (Branch \code{0\_mortarcontact})}
\author{\textbf{Manuel Hradsky}\\
\small Arbeitsbereich f\"ur Festigkeitslehre und Baustatik\\
\small Universit\"at Innsbruck}
\date{\today}

\begin{document}

\maketitle

\begin{abstract}
Dieses Dokument beschreibt Theorie, Herleitung und Implementierung des \emph{reibungsfreien}
Mortar-Normalkontakts, wie er auf dem Branch \code{0\_mortarcontact} des Finite-Elemente-Frameworks
\code{EdelweissFE} rechnet. Die Darstellung ist bewusst als in sich geschlossene, paper-artige
Referenz angelegt: jede Formel, jeder Algorithmus und jede Konstante wird aus den Primärquellen
der Literatur hergeleitet und mit exakter Fundstelle (Autor, Jahr, Gleichung bzw.\ Abschnitt)
belegt, und jeder Baustein wird auf die konkrete Codestelle (\file{Datei:Zeile}) abgebildet.
Die Formulierung beruht auf \emph{dualen} Lagrange-Multiplikatoren zur knotenweisen Entkopplung der
Kontaktbedingungen \citep{wohlmuth2000}, einer \emph{segmentbasierten} Integration mit linearer
Sub-Zellen-Zerlegung für lineare und quadratische Facetten- und Linienelemente
(\code{CONLINE2/3}, \code{CONQUAD4/8/9}, \code{CONTRI3/6}), einer konsistenten Randbehandlung
partiell überdeckter Elemente \citep{cichosz} und einer semiglatten Primal-Dual
Active-Set-Strategie (PDASS) zur Lösung der diskreten Signorini-Bedingungen
\citep{hueber2005,gitterle}. Besonderes Gewicht liegt auf den vier
technischen Kernpunkten der zugehörigen Präsentation: (i)~warum die Segmentquadratur den
Genauigkeitsgrad~5 verwendet (7-Punkt-Dreiecksregel nach \citealt{dunavant1985}), (ii)~Herleitung der
Koppelmatrizen $\bm D,\bm M$ und des gewichteten Spalts, (iii)~warum die Basistransformation
quadratischer Elemente den Faktor $\tfrac13$ trägt \citep{popp2012}, und (iv)~die Herleitung
der Normal-NCP-Funktion \citep[Gl.~(55)]{gitterle}. Reibung und statische Kondensation
sind auf diesem Branch \emph{nicht} enthalten und werden hier bewusst nicht behandelt.
\end{abstract}

\clearpage
\tableofcontents
\clearpage

% =====================================================================
\section*{Geltungsbereich und Implementierungsstand}
\addcontentsline{toc}{section}{Geltungsbereich und Implementierungsstand}

Beschrieben wird ausschließlich der Zustand auf Branch \code{0\_mortarcontact}: der reibungsfreie
Mortar-Normalkontakt in \textbf{Sattelpunkt-Formulierung} (die Lagrange-Multiplikatoren sind
explizite zusätzliche Unbekannte). Es gibt auf diesem Branch:
\begin{itemize}
  \item \textbf{keine} Coulomb-Reibung (kein \code{friction\_coefficient}-Schlüsselwort),
  \item \textbf{keine} statische/duale Kondensation (kein \code{condensation}-Schlüsselwort). Die
        Constraint-Argumente sind \code{nonMortarSurface}, \code{mortarSurface}, \code{field} sowie
        der NCP-Parameter \code{cn} (\file{mortarcontact.py:\allowbreak{}89--108}); \code{cn} ist optional
        (Standard $10^6$), sollte aber gesetzt werden (Abschnitt~\ref{sec:ncp}).
\end{itemize}
Auf diesem Branch wird pro Slave-Knoten \emph{genau ein skalarer} Lagrange-Multiplikator
instanziiert -- die Normalkomponente $\lambda_I$ der Kontakttraktion (der Kontaktdruck): es gibt
einen Multiplikator-Freiheitsgrad je Slave-Knoten, \emph{nicht} $\dim$
(\file{mortarcontact.py:\allowbreak{}348--351}). Die Kontaktkraft entsteht daraus als
$\lambda_I\,\bm n_I$, also allein in Richtung der Knotennormale
(\file{mortarcontact.py:\allowbreak{}825, 833}). In der allgemeinen (kontinuierlichen)
Mortar-Formulierung ist der Multiplikator ein vollst\"andiger Traktionsvektor
$\bm z_I\in\mathbb R^{\dim}$, dessen Tangentialkomponenten erst mit Coulomb-Reibung physikalisch
aktiv w\"urden (dann $\dim$ Multiplikatoren je Knoten). Diese Vektor-Erweiterung ist der
nat\"urliche Reibungspfad, auf Branch \code{0\_mortarcontact} aber \emph{nicht} instanziiert.
Der Code ist reines Python im Framework \code{EdelweissFE}; die drei
zentralen Quelldateien -- \file{mortarcontact.py} (867~Z.), \file{mortar\_geom\_utils.py} (119~Z.)
und \file{contactelement/element.py} (511~Z.) -- werden mit ihren vollständigen Pfaden in
Abschnitt~\ref{sec:architektur} (Tab.~Modulzuordnung) aufgeführt.

\medskip
\noindent\textbf{Notationskonvention.} In der Literatur und in der Präsentation heißt die
Slave--Master-Koppelmatrix~$\bm M$. Im Quelltext heißt dieselbe Matrix~\code{C}, während der
Bezeichner~\code{M} dort für die lokale Biorthogonalitäts-Massenmatrix reserviert ist. Dieses
Dokument folgt der Literatur und schreibt durchgehend $\bm D$ (Slave--Slave) und $\bm M$
(Slave--Master); an den Codestellen wird jeweils auf den Bezeichner \code{C} hingewiesen.

% =====================================================================
\section{Einleitung und Problemstellung}
\label{sec:einleitung}

Zwei deformierbare Körper -- im Zielanwendungsfall ein Stahlanker (Master) in einem Betonblock
(Slave) -- sollen an ihrer gemeinsamen Kontaktfläche $\gamma_c$ Kräfte übertragen, ohne sich zu
durchdringen. Die Finite-Elemente-Netze der beiden Körper sind \emph{nichtkonform}: ihre Knoten
treffen an der Kontaktfläche nicht aufeinander. Ein klassischer Knoten-auf-Knoten-Kontakt
(\emph{node-to-node}) setzt die Nichtdurchdringung punktweise an paarenden Knoten durch und ist
damit auf nichtkonforme Netze nicht anwendbar: es gibt keine natürliche Knotenpaarung, und ein
naiver Knoten-auf-Segment-Kontakt (\emph{node-to-segment}) verletzt den Kontakt-Patch-Test, weil er
einen konstanten Druck nicht exakt überträgt~\citep{wriggers}.

Die \textbf{Mortar-Methode} setzt die Kontaktbedingung stattdessen \emph{schwach} durch: nicht
punktweise, sondern im integralen Mittel über die Kontaktfläche, gewichtet mit den Formfunktionen
der Slave-Seite. Der zugehörige Lagrange-Multiplikator $\bm\lambda$ ist physikalisch die
Kontakttraktion (im reibungsfreien Fall der Kontaktdruck) und lebt auf der Slave-Fläche. Werden für
$\bm\lambda$ \emph{duale} Formfunktionen im Sinne von \citet{wohlmuth2000} gewählt, so
entkoppeln die Kontaktbedingungen knotenweise (Abschnitte~\ref{sec:koppelmatrizen}
und~\ref{sec:dual}), und die Methode besteht den Kontakt-Patch-Test bei optimaler
Konvergenzordnung. Die Mortar-Methode ist heute der Standardzugang für Kontakt zwischen
nichtkonformen Netzen bei großen Deformationen~\citep{puso2004,popp2010,farah}.

Dieses Dokument beschreibt die konkrete reibungsfreie Umsetzung in \code{EdelweissFE}. Die Darstellung
folgt der Verarbeitungs-Pipeline: von der Geometrie (Nachbarsuche, Projektion, Clipping,
Quadratur) über die diskreten Koppelmatrizen und den gewichteten Spalt bis zur Lösung der
Kontaktbedingung als semiglattes Komplementaritätsproblem.

% =====================================================================
\section{Kontinuierliche Formulierung}
\label{sec:kontinuierlich}

\subsection{Kinematik und Spalt}
Der Slave-Körper trägt den Index~$(1)$, der Master-Körper den Index~$(2)$. Zu einem Slave-Punkt
$\bm x^{(1)}\in\gamma_c$ wird über die Kontaktnormale $\bm n$ der nächste Master-Punkt
$\hat{\bm x}^{(2)}$ (die \emph{Projektion}) bestimmt. Der Normalspalt ist
\begin{equation}
g_n = -\,\bm n\cdot\big(\bm x^{(1)} - \hat{\bm x}^{(2)}\big),
\label{eq:gap}
\end{equation}
mit der Vorzeichenkonvention $g_n>0$ bei geöffnetem Kontakt (Trennung) und $g_n<0$ bei
Durchdringung~\citep{wriggers,farah}.

\subsection{Hertz--Signorini--Moreau / KKT-Bedingungen}
Der reibungsfreie Normalkontakt fordert Nichtdurchdringung, nichtklebende (nur drückende)
Kontakttraktion und Komplementarität. Mit dem Kontaktdruck $p=\lambda\ge0$ (Normalkomponente der
Traktion) lauten die Karush--Kuhn--Tucker-Bedingungen (Hertz--Signorini--Moreau)
\begin{equation}
g_n \ge 0, \qquad p \ge 0, \qquad p\, g_n = 0 \qquad \text{auf } \gamma_c .
\label{eq:kkt}
\end{equation}
Entweder ist der Spalt offen ($g_n>0$) und der Druck null, oder es besteht Kontakt ($g_n=0$) mit
$p\ge0$~\citep[Kap.~4]{wriggers}.

\subsection{Schwache Form und Multiplikatorraum}
Die Kontaktbedingung wird über einen Lagrange-Multiplikator $\bm\lambda$ (die Kontakttraktion) in
das Prinzip der virtuellen Arbeit eingebracht. Der zusätzliche Kontaktbeitrag zur virtuellen Arbeit
ist
\begin{equation}
\delta W_c = \int_{\gamma_c} \bm\lambda\cdot\big(\delta\bm u^{(1)} - \delta\bm u^{(2)}\big)\,
\mathrm d\gamma,
\end{equation}
und die schwache Nichtdurchdringung wird mit Testfunktionen $\delta\bm\lambda$ aus dem
Multiplikatorraum $M$ als Variationsungleichung angesetzt. Die Wahl des diskreten Raums für
$\bm\lambda$ ist entscheidend: \citet{wohlmuth2000} zeigt, dass ein \emph{dualer} Raum
(biorthogonal zu den Slave-Formfunktionen) die inf-sup/LBB-Stabilität sichert und eine optimale
a-priori-Fehlerabschätzung liefert, während er zugleich die spätere knotenweise Elimination erlaubt.
Genau dieser duale Raum wird hier verwendet (Abschnitt~\ref{sec:dual}). Die Diskretisierung von
\eqref{eq:gap}--\eqref{eq:kkt} führt auf die gewichteten Größen der folgenden Abschnitte: die
Koppelmatrizen $\bm D,\bm M$, den gewichteten Spalt $\tilde g_I$ und die knotenweise
NCP-Bedingung.

% =====================================================================
\section{Diskretisierung und Programmarchitektur}
\label{sec:architektur}

\subsection{Kontaktelemente als geometrische Rand-Overlays}
\code{EdelweissFE} besitzt keine eingebauten Kontaktelemente. Der Mortar-Kontakt definiert sie daher
selbst als rein \emph{geometrische} Rand-Overlay-Elemente (kein Materialkern): sie tragen die
Verschiebungs-Freiheitsgrade der zugehörigen Volumenoberfläche und liefern Formfunktionen $N_I$,
duale Funktionen $\Phi_I$ und Quadraturpunkte. Erzeugt werden sie aus den Seitenflächen (Sidesets)
der Volumenelemente (z.\,B.\ CONQUAD8 auf der Oberfläche eines hex20).

\begin{center}
\small
\begin{tabular}{lccc}
\toprule
Typ & Knoten & Geometrie & Ordnung \\
\midrule
CONLINE2 / CONLINE3 & 2 / 3 & 2D-Kante & linear / quadratisch \\
CONQUAD4 / CONQUAD8 / CONQUAD9 & 4 / 8 / 9 & 3D-Fläche & linear / quadratisch \\
CONTRI3 / CONTRI6 & 3 / 6 & 3D-Dreieck & linear / quadratisch \\
\bottomrule
\end{tabular}
\end{center}

\subsection{Verarbeitungs-Pipeline}
Die Berechnung folgt einer festen Pipeline, die zugleich die Gliederung dieses Dokuments ist:
\begin{enumerate}\itemsep2pt
  \item \textbf{Nachbarsuche} (BVH) $\to$ Kandidatenpaare (Abschnitt~\ref{sec:geometrie}).
  \item \textbf{Projektion + Clipping + Triangulierung} $\to$ Integrationszellen
        (Abschnitt~\ref{sec:geometrie}), für quadratische Facetten nach linearer Sub-Cell-Zerlegung
        (Abschnitt~\ref{sec:subcells}).
  \item \textbf{Segmentquadratur} (Grad~5) $\to$ Koppelmatrizen $\bm D,\bm M$
        (Abschnitte~\ref{sec:quadratur}, \ref{sec:koppelmatrizen}).
  \item \textbf{Gewichteter Spalt} $\tilde g_j$ und \textbf{NCP} $\to$ Active Set
        (Abschnitte~\ref{sec:koppelmatrizen}, \ref{sec:ncp}).
  \item \textbf{Assemblierung + Newton} $\to$ Lösung (Abschnitt~\ref{sec:solve}).
\end{enumerate}

\subsection{Modulzuordnung}
\begin{center}
\small
\begin{tabularx}{\linewidth}{lX}
\toprule
Datei & Rolle \\
\midrule
\file{constraints/mortarcontact.py} & Segmentierung, $\bm D,\bm M$-Assemblierung, gewichteter Spalt, NCP/Active-Set, K/PExt-Assemblierung \\
\file{constraints/mortar\_geom\_utils.py} & Ebenenprojektion, Sutherland--Hodgman-Clipping, Triangulierung, 1D-Clip \\
\file{elements/contactelement/element.py} & Formfunktionen, Quadraturpunkte, Basistransformation $\bm T_e$, lokale $\bm M_e/\bm D_e/\bm A_e$ \\
\bottomrule
\end{tabularx}
\end{center}

% =====================================================================
\section{Geometrische Verarbeitung: Nachbarsuche, Projektion, Clipping}
\label{sec:geometrie}

Bei nichtkonformen Netzen ist a priori unbekannt, welche Master-Facette welche Slave-Facette
überlappt. Die geometrische Verarbeitung bestimmt diese Überlappungen und zerlegt sie in
Integrationszellen (Folien~5--6).

\subsection{Nachbarsuche mit Bounding-Volume-Hierarchie}
Eine naive Prüfung aller Slave--Master-Paare kostet $\mathcal O(N^2)$. Stattdessen werden die
achsparallelen Hüllboxen (AABB) der Master-Facetten in einer \emph{Bounding-Volume-Hierarchie}
organisiert -- einem binären Baum mit Median-Split entlang der längsten Achse und Blattgröße~$\le2$
\citep{ericson2004}. Pro Slave-Facette liefert eine Baumabfrage nur wenige
\emph{Kandidaten}; der Aufwand sinkt auf $\sim\mathcal O(N\log N)$. Die AABB erhalten einen
Sicherheitsrand $\max(0.15\cdot\text{Ausdehnung},\,0.1)$; nur für die Kandidaten wird anschließend
die exakte geometrische Verschneidung durchgeführt.

\subsection{Hilfsebene und Projektion}
Für ein Kandidatenpaar wird an der Slave-(Sub-)Facette eine \emph{Hilfsebene} definiert (Stützpunkt
$\bm p_0$, Normale $\bm n$). Beide Polygone -- Slave und Master -- werden in diese Ebene projiziert
und dort in ebenen Koordinaten dargestellt (\file{mortar\_geom\_utils.py}, \code{get\_tangent\_basis},
\code{to\_plane\_coords}). Diese Ebene dient nur der Schnittbildung; die tatsächliche Geometrie geht
später über die Newton-Rückabbildung an den Gauß-Punkten ein (Abschnitt~\ref{sec:quadratur}).

\subsection{Sutherland--Hodgman-Clipping und Triangulierung}
In der Ebene wird das Master-Polygon gegen das (konvexe) Slave-Polygon geschnitten. Dafür dient der
\emph{Sutherland--Hodgman-Algorithmus}~\citep{sutherland1974}: das zu schneidende Polygon wird
sukzessive an jeder Kante des Clip-Polygons gekappt (\emph{reentrant polygon clipping}). Der
Algorithmus setzt ein \emph{konvexes} Clip-Polygon voraus -- deshalb wird als Clip-Fenster die
(konvexe) Slave-Facette bzw.\ Sub-Cell verwendet, nicht umgekehrt. Das Ergebnis ist das konvexe
Überlappungspolygon $\text{Slave}\cap\text{Master}$ (\code{sutherland\_hodgman\_clip}).

Da die Segmentquadratur (Abschnitt~\ref{sec:quadratur}) auf Dreiecken definiert ist, wird das
Überlappungspolygon abschließend \emph{fächerartig} von einer Ecke aus trianguliert
(\code{triangulate\_polygon}): ein $n$-Eck ergibt $n-2$ Integrationsdreiecke. Weil das Polygon konvex
ist, ist diese Triangle-Fan-Zerlegung immer gültig. Im 2D-Fall reduziert sich das Ganze auf einen
1D-Intervallschnitt (\code{clip\_1d\_segments}). Das Vorgehen entspricht MOOSEs
\code{AutomaticMortarGeneration}.

% =====================================================================
\section{Segmentbasierte Quadratur: warum Genauigkeitsgrad 5}
\label{sec:quadratur}

Nach dem Clipping (Abschnitt~\ref{sec:geometrie}) liegt die Überlappung eines Slave--Master-Paars
als Menge ebener Integrationsdreiecke (3D) bzw.\ Teilstrecken (2D) vor. Auf diesen wird integriert.
Dieser Abschnitt beantwortet die zentrale Frage der Präsentation (Folie~7): \emph{Warum verwendet
der Code Genauigkeitsgrad~5 -- konkret die 7-Punkt-Dreiecksregel im 3D-Fall?} Die Antwort hat drei
Schichten: (i)~warum überhaupt \emph{segment}basiert integriert wird, (ii)~warum der effektive
Integrationsgrad \emph{über} das nominelle Formfunktionsprodukt angehoben werden muss, und
(iii)~woher die konkrete 7-Punkt-Regel stammt.

\subsection{Notwendigkeit der segmentbasierten Integration}
Die Mortar-Matrizen $\bm D,\bm M$ (Abschnitt~\ref{sec:koppelmatrizen}) enthalten unter dem Integral
ein Produkt aus Größen \emph{beider} Seiten -- der dualen bzw.\ Standard-Formfunktionen der
Slave-Seite und den auf die Slave-Fläche zurückprojizierten Formfunktionen der Master-Seite. Bei
nichtkonformen Netzen besitzt dieser Integrand \emph{Knicke} (schwache Diskontinuitäten) im Inneren
einer Slave-Facette: nämlich überall dort, wo eine Master-Elementkante die Slave-Facette kreuzt.
\citet{farah2015} formulieren das explizit:
\begin{quote}
\small ``the integrand represents a non-smooth function which cannot be evaluated exactly by using
standard Gauss rules. This non-smoothness stems from the locally supported Lagrange polynomials on
master and slave side having kinks at the respective element nodes and edges.''
\hfill \citep[S.~214]{farah2015}
\end{quote}
Eine Standard-Gauß-Regel, die über die ganze Facette gelegt wird (\emph{element}based), ``sieht''
diese Knicke nicht und ist deshalb hochgradig empfindlich gegenüber Gauß-Punkt-Zahl und
Netz-Verhältnis~\citep[S.~217]{farah2015}. Der Ausweg ist die \emph{segment}basierte Integration:
die Überlappung wird in glatte, knickfreie Zellen mit konstantem Jacobian zerlegt (hier: Dreiecke),
auf denen der Integrand wieder polynomial-glatt ist und Gauß-Quadratur ihre volle Ordnung
entfaltet~\citep[Abschn.~3.1]{farah2015}. Genau das leistet die Clipping--Triangulierungs-Pipeline aus
Abschnitt~\ref{sec:geometrie}.

\subsection{Anhebung des Integrationsgrades}
Innerhalb einer knickfreien Integrationszelle bleibt ein Resteffekt: die Rückabbildung eines
Gauß-Punkts von der Hilfsebene in die \emph{Natur}koordinaten von Slave und Master ist
\emph{nichtlinear} (bei gekrümmten/verzerrten Facetten und generell durch die Projektion). Dadurch
ist der Integrand -- als Funktion der Zell-Koordinaten -- kein Polynom vom Grad des reinen
Formfunktionsprodukts, sondern effektiv höhergradig. \citet{farah2015} geben deshalb die
Faustregel und die konkrete Empfehlung an (der Satz findet sich wortgleich in \citealt[Anh.~A.1.1]{farah}):
\begin{quote}
\small ``the highest polynomial degree that can be computed exactly by the integration rule should
be chosen somewhat higher than the product of shape functions on slave and master side would
indicate. This is due to the nonlinearity of the projections between auxiliary plane and actual
contact surfaces. Based on our experiences, we suggest to use seven Gauss points per triangular
integration cell, which are able to exactly calculate a polynomial degree of 5.''
\hfill (\citealp[S.~217]{farah2015}; identisch \citealp[Anh.~A.1.1, S.~209--210]{farah})
\end{quote}
Die Wahl ``7 Punkte / Grad~5'' ist also eine \emph{empirisch abgesicherte} Empfehlung: sie liegt im
gesättigten Genauigkeitsbereich, höhere Punktzahlen ändern die Lösung praktisch nicht mehr
(\citealp[Kap.~6]{farah}). Im Code ist genau das umgesetzt (\file{mortarcontact.py:\allowbreak{}201--224}), mit
demselben Begründungskommentar und Verweis auf \citet[Anh.~A.1.1]{farah}.

\subsection{Die 7-Punkt-Regel vom Grad 5 (Dunavant, 1985)}
Die symmetrische 7-Punkt-Regel mit Genauigkeitsgrad~5 auf dem Referenzdreieck stammt aus
\citet{dunavant1985}. Er approximiert (Gl.~(4), S.~1129)
\begin{equation}
\int_A f(\alpha,\beta,\gamma)\,\mathrm dA \;=\; A\sum_{i=1}^{n_g} w_i\, f(\alpha_i,\beta_i,\gamma_i),
\qquad \alpha+\beta+\gamma=1,
\end{equation}
in baryzentrischen Koordinaten und bestimmt die $(w_i,\alpha_i,\beta_i,\gamma_i)$ aus den
Momentenbedingungen $\int_A \alpha^i\beta^j\,\mathrm dA = 2A\,i!\,j!/(i+j+2)!$ (Gl.~(13), S.~1132).
Die Punkte werden in \emph{Symmetrieorbits} gruppiert: ein Schwerpunktpunkt ($n_0$), Dreiergruppen
auf den Seitenhalbierenden ($n_1$, je 3~Punkte) und allgemeine Sechsergruppen ($n_2$), mit
Gesamtzahl $n_g = n_0 + 3n_1 + 6n_2$ (Gl.~(34), S.~1135). Für Grad $p=5$ liefert Dunavant
$n_0=1,\;n_1=2,\;n_2=0$, also
\begin{equation}
n_g = 1 + 3\cdot 2 = 7 \quad\text{(Dunavant, 1985, Tab.~III/IV, S.~1135--1137)}.
\end{equation}
Das ist effizienter als die Gauß-Produktregel, die für Grad~5 neun Punkte bräuchte
($7<9$). Die konkreten Werte \citep[Anh.~II, S.~1140]{dunavant1985} sind:
\begin{center}
\small
\begin{tabular}{lcccc}
\toprule
Orbit & Gewicht $w_i$ (pro Punkt) & $\alpha$ & $\beta,\gamma$ & \#Punkte \\
\midrule
Schwerpunkt & $0.225000000000$ & $0.333333333333$ & $0.333333333333$ & 1 \\
$n_1$-Orbit 1 & $0.132394152789$ & $0.059715871790$ & $0.470142064105$ & 3 \\
$n_1$-Orbit 2 & $0.125939180545$ & $0.797426985353$ & $0.101286507323$ & 3 \\
\bottomrule
\end{tabular}
\end{center}
Die Dreiergruppen entstehen durch zyklische Permutation von $(\alpha,\beta,\gamma)$; die
Gewichtssumme ist $0.225 + 3\cdot0.132394\ldots + 3\cdot0.125939\ldots = 1.000$. Zwei Eigenschaften
machen gerade diese Regel für den Kontakt attraktiv: \textbf{alle sieben Punkte liegen im Inneren}
des Dreiecks ($\alpha,\beta,\gamma\in(0,1)$) und \textbf{alle Gewichte sind positiv}. Grad~5 ist
oberhalb von Grad~2 die niedrigste Dunavant-Regel mit diesen beiden Eigenschaften -- die Grad-3-Regel
etwa hat ein negatives Gewicht ($-0.5625$; \citealp[S.~1140]{dunavant1985}) und ist damit numerisch
ungünstig. Dieselbe Regel steht im Code in \file{mortarcontact.py:\allowbreak{}201--224}; die
Gauß-Punkte sind identisch, die Gewichte dort ($9/80,\ (155\mp\sqrt{15})/2400$, Summe~$\tfrac12$
statt~$1$) sind bereits mit der Referenzdreiecksfläche~$\tfrac12$ verrechnet, sodass das Produkt
$w\cdot J_{\text{cell}}$ und damit die Regel unverändert bleibt.

\subsection{Der zweidimensionale Fall}
Im 2D-Fall wird pro geclipptem Liniensegment eine 3-Punkt-Gauß--Legendre-Regel auf $[-1,1]$ mit
Punkten $\xi\in\{0,\pm\sqrt{0.6}\}$ und Gewichten $\{5/9,8/9,5/9\}$ eingesetzt
(\file{mortarcontact.py:\allowbreak{}192--198}). Sie ist exakt bis Polynomgrad $2n-1=5$ (Nullstellen des
Legendre-Polynoms $P_3$) -- daher der Kommentar ``degree~5'' im Code, konsistent mit der Grad-5-Wahl
im 3D-Fall. \emph{Bemerkung:} \citet{farah} verwendet im 2D-Kontext empirisch \emph{fünf}
Gauß-Punkte pro Segment, nicht drei. Beide Wahlen folgen derselben Logik (Grad über das
Formfunktionsprodukt anheben wegen nichtlinearer Projektion); der Code priorisiert hier die
einheitliche \emph{Grad}-5-Aussage über 2D und 3D. Für die im 2D auftretenden Integranden ist Grad~5
in der Praxis ausreichend (bestätigt durch die 2D-Patch-Tests, Abschnitt~\ref{sec:verif}).

\subsection{Auswertung an den Gauß-Punkten}
An jedem Gauß-Punkt wird der Ebenenpunkt per \emph{Newton-Iteration} in die Naturkoordinaten von
Slave \emph{und} Master zurückabgebildet, sodass dort $N_s$ und $N_m$ (und die dualen $\Phi$)
ausgewertet werden können; das gewichtete Produkt geht mit $w\cdot J_{\text{cell}}$ in die
Summation über alle Zellen ein und liefert $\bm D$ und $\bm M$ (\file{mortarcontact.py:\allowbreak{}600--648};
Summenform nach \citealt[Alg.~1, S.~215]{farah2015}):
\begin{equation}
\bm D_{IK} = \sum_{c=1}^{n_{\text{cell}}}\sum_{g=1}^{n_g} w_g\,
\Phi_I(\xi_g^{(1)})\,N_K^{(1)}(\xi_g^{(1)})\,J_c,
\qquad
\bm M_{IJ} = \sum_{c}\sum_{g} w_g\,
\Phi_I(\xi_g^{(1)})\,N_J^{(2)}(\xi_g^{(2)})\,J_c .
\end{equation}

% =====================================================================
\section{Koppelmatrizen $\bm D,\bm M$ und gewichteter Spalt}
\label{sec:koppelmatrizen}

Dieser Abschnitt leitet die diskreten Mortar-Koppelmatrizen und den gewichteten Spalt her -- die
Größen der Präsentation (Folie~8), die am Ende der Segmentquadratur stehen.

\subsection{Entstehung von $\bm D$ und $\bm M$}
Man diskretisiert die Kontakttraktion mit dem dualen Multiplikatorfeld
$\bm\lambda_h = \sum_j \Phi_j\,\bm z_j$ (allgemeine Vektorform; auf Branch \code{0\_mortarcontact}
wird davon nur die Normalkomponente $\bm z_j\!\cdot\!\bm n_j=\lambda_j$ instanziiert, s.\
Geltungsbereich) und die Geometrie beider Seiten mit den jeweiligen
Standard-Formfunktionen $N^{(1)},N^{(2)}$ und setzt beides in die Kontakt-Virtuelle-Arbeit ein.
Sortiert man nach den Knotenbeiträgen um, treten genau zwei Matrizen auf -- die \emph{Mortar-Matrizen}
(\citealp[Gl.~(3.34)--(3.36)]{farah}; \citealp[Gl.~(3.5)--(3.6)]{popp2012}; \citealp[Gl.~(35)/(36)]{poppgeewall2009}):
\begin{equation}
\bm D[j,k] = D_{jk}\,\bm I \;=\; \int_{\gamma_{c,h}^{(1)}} \Phi_j\,N_k^{(1)}\,\mathrm d\gamma\;\bm I,
\qquad
\bm M[j,l] = M_{jl}\,\bm I \;=\; \int_{\gamma_{c,h}^{(1)}} \Phi_j\,\big(N_l^{(2)}\!\circ P_h\big)\,\mathrm d\gamma\;\bm I .
\label{eq:DM}
\end{equation}
Dabei ist $j$ ein Slave-Multiplikatorknoten, $k$ ein Slave-Verschiebungsknoten $(1)$, $l$ ein
Master-Verschiebungsknoten $(2)$, $\bm I\in\mathbb R^{\dim\times\dim}$ die Einheitsmatrix, und
$P_h:\gamma_{c,h}^{(1)}\!\to\!\gamma_{c,h}^{(2)}$ die diskrete Kontaktabbildung (die Projektion). Beide
Integrale laufen über die \emph{Slave}-Fläche. $\bm D$ (Slave--Slave) ist quadratisch, $\bm M$
(Slave--Master) im Allgemeinen rechteckig~\citep[S.~B427]{popp2012}. Die algebraische Kontaktkraft ist
$\bm f_c = [\,\bm 0\;\;-\bm M\;\;\bm D\,]^{\!\top}\bm z$ \citep[Gl.~(3.38)]{farah}. Im Code werden
$\bm D,\bm M$ in \file{mortarcontact.py:\allowbreak{}487--690} assembliert (die Slave--Master-Matrix heißt dort
\code{C}; vgl.\ Notationskonvention).

\subsection{Gewichteter (schwacher) Normalspalt}
Die diskrete schwache Nichtdurchdringung wird pro Slave-Knoten $j$ zu einem \emph{gewichteten Spalt}
\citep[Gl.~(3.40)]{farah}
\begin{equation}
\tilde g_{n,j} = \int_{\gamma_{c,h}^{(1)}} \Phi_j\, g_{n,h}\,\mathrm d\gamma .
\label{eq:wgap_int}
\end{equation}
Setzt man $g_{n,h}$ über \eqref{eq:gap} ein und sortiert nach Knotenpositionen um, erhält man die
Positionsform (\citealp[Gl.~(36) und~(50)]{gitterle}; \citealp[Gl.~(50)]{poppgeewall2009}):
\begin{equation}
\tilde g_{j} \;=\; \bm n_j\!\cdot\!\Big(\sum_k \bm D[j,k]\,\bm x_k^{(1)}
- \sum_l \bm M[j,l]\,\bm x_l^{(2)}\Big)
\;=\; \bm n_j\!\cdot\!\Big(\bm D[j,j]\,\bm x_j^{(1)} - \sum_l \bm M[j,l]\,\bm x_l^{(2)}\Big),
\label{eq:wgap_pos}
\end{equation}
wobei die zweite Gleichheit gilt, weil $\bm D$ bei dualer Basis diagonal wird
(Abschnitt~\ref{sec:dual}). Genau diese Form steht im Code (\file{mortarcontact.py:\allowbreak{}785--795},
Bezeichner \code{C} für $\bm M$).

\medskip
\noindent\textbf{Warum schwach/nodal statt punktweise?} Der gewichtete Spalt ist ein volumetrisches
(integrales) Maß, das dennoch als knotenweises Näherungsmaß verwendet wird
(\citealp{farah}: \emph{``The discrete weighted gap $\tilde g_{n,j}$ represents a volumetric measure
\ldots\ Nevertheless, it is utilized as node-wise measure''}). Der tiefere Grund für die schwache
Auswertung liegt in der Ungleichungsnatur der Nichtdurchdringung in Kombination mit
Formfunktionen höherer Ordnung: eine punktweise (starke) Durchsetzung ist damit nicht konsistent
behandelbar~\cite[S.~B429]{popp2012}. Die knotenweisen HSM-Bedingungen lauten dann
$\tilde g_{n,j}\ge0,\; \lambda_{n,j}\ge0,\; \lambda_{n,j}\tilde g_{n,j}=0$
(\citealp[Gl.~(3.7)]{popp2012}; \citealp[Gl.~(3.41)]{farah}) -- die diskrete Fassung von \eqref{eq:kkt}.

\subsection{Zeilensummen-Identität und Impulserhaltung}
Eine zentrale Konsistenzeigenschaft ist die Zeilensummen-Identität
\begin{equation}
\sum_k \bm D[j,k] \;=\; \sum_l \bm M[j,l] \qquad \forall\, j
\label{eq:rowsum}
\end{equation}
\citep[Gl.~(4.99)]{farah}. Sie folgt aus der Erhaltung des linearen Impulses
$\bm f_\lambda^{(1)}-\bm f_\lambda^{(2)}=\bm 0$: mit der Partition-of-Unity der Formfunktionen
$\sum_k N_k^{(1)}=\sum_l N_l^{(2)}=1$ \citep[Gl.~(4.97)]{farah} reduziert sich die Differenz der
Zeilensummen auf $\int\Phi_j\,\mathrm d\gamma - \int\Phi_j\,\mathrm d\gamma = 0$, was
\emph{immer} erfüllt ist, sofern $\bm D$ und $\bm M$ über \emph{dasselbe} diskrete Gebiet integriert
werden (\citealp[Gl.~(4.100)--(4.101)]{farah}; identisch \citealt{popp2010} und \citealt{puso2004}). Physikalisch
sichert \eqref{eq:rowsum} die Translationsinvarianz des Spalts (eine starre Translation beider Körper
erzeugt keinen künstlichen Gap) und ist Voraussetzung für das Bestehen des Kontakt-Patch-Tests. Genau
deshalb müssen $\bm D$ und $\bm M$ im Code über \emph{dieselben} Integrationszellen summiert werden
(die konsistente Randbehandlung, Abschnitt~\ref{sec:subcells}, stellt das auch für partiell
überdeckte Elemente sicher).

% =====================================================================
\section{Duale Lagrange-Multiplikatoren und Biorthogonalität}
\label{sec:dual}

\subsection{Definition und Motivation}
Die dualen Formfunktionen $\Phi_j$ sind durch die \emph{Biorthogonalität} zu den
Slave-Formfunktionen definiert (\citealp[Gl.~(3.5)]{wohlmuth2000}; \citealp[Gl.~(4.1)]{popp2012};
\citealp[Gl.~(30)]{poppgeewall2009}; \citealp[Gl.~(3.46)]{farah}):
\begin{equation}
\int_{\gamma_{c,h}^{(1)}} \Phi_j\, N_k^{(1)}\,\mathrm d\gamma
\;=\; \delta_{jk}\int_{\gamma_{c,h}^{(1)}} N_k^{(1)}\,\mathrm d\gamma .
\label{eq:biorth}
\end{equation}
\citet{wohlmuth2000} führte diese dualen Räume ein, um bei gleicher \emph{Optimalität}
(inf-sup/LBB-Stabilität mit $h$-unabhängiger Konstante, optimale a-priori-Fehlerabschätzung der
Ordnung $\mathcal O(h)$ in der Energienorm) lokal getragene Multiplikator-Basisfunktionen zu erhalten:
\begin{quote}
\small ``this Lagrange multiplier space is replaced by a dual space without losing the optimality of
the method \ldots\ all the basis functions of the new method are supported in a few elements. The
mortar map can be represented by a diagonal matrix; in the standard mortar method a linear system of
equations must be solved.'' \hfill \citep[Abstract]{wohlmuth2000}
\end{quote}

\subsection{Diagonalisierung von $\bm D$}
Setzt man \eqref{eq:biorth} in die Definition~\eqref{eq:DM} von $\bm D$ ein, verschwinden alle
Außerdiagonalterme:
\begin{equation}
D_{jk} = \int \Phi_j N_k^{(1)}\,\mathrm d\gamma = \delta_{jk}\int N_k^{(1)}\,\mathrm d\gamma,
\qquad\Longrightarrow\qquad \bm D = \operatorname{diag}\!\Big(\textstyle\int N_k^{(1)}\,\mathrm d\gamma\Big).
\end{equation}
Die gesamte algebraische Struktur nimmt damit die Gestalt eines Knoten-auf-Segment-Kontakts an, und
$\bm D$ ist trivial invertierbar (\citealp[Abschn.~3.4.1.2, S.~33]{farah}: \emph{``the algebraic form of the slave
side mortar matrix D \ldots\ becomes a diagonal matrix''}). Auf diese Diagonalität stützt sich der
Projektionsoperator $\bm P = \bm D^{-1}\bm M$ \citep[Gl.~(3.65)]{farah} -- die algebraische Grundlage
jeder späteren Kondensation. \emph{Anmerkung:} Diese statische Kondensation ist auf Branch
\code{0\_mortarcontact} bewusst \emph{nicht} aktiviert (Sattelpunkt-Formulierung); die Diagonalität
von $\tilde{\bm D}$ wird hier nur für die Normierung im NCP-Indikator (Abschnitt~\ref{sec:ncp}) genutzt.

\subsection{Elementlokale Konstruktion}
Bei Elementen mit nicht-konstanter Jacobi-Determinante werden die $\Phi_j$ als Linearkombination
der Standard-Formfunktionen dargestellt, $\Phi_j = c_{jk}N_k$, mit der Koeffizientenmatrix
(\citealp[Gl.~(3.53)--(3.56)]{farah})
\begin{equation}
\bm A_e \;=\; \bm D_e\,\bm M_e^{-1},
\qquad
(\bm D_e)_{jk} = \delta_{jk}\!\int_e N_k\,J\,\mathrm de,
\qquad
(\bm M_e)_{jk} = \int_e N_j N_k\,J\,\mathrm de .
\label{eq:Ae}
\end{equation}
Hier sind $\bm D_e,\bm M_e$ \emph{lokale Hilfs-Massenmatrizen} zur Konstruktion der dualen Basis --
nicht zu verwechseln mit den globalen Koppelmatrizen $\bm D,\bm M$ (Farah unterscheidet sie im
Symbolverzeichnis). Im Code ist \eqref{eq:Ae} zweifach umgesetzt: als Referenz-Fallback in
\file{element.py:\allowbreak{}470--508} (\code{A\_e = D\_e @ inv(M\_e) @ T\_e}) und -- bevorzugt -- zur Laufzeit
über die \emph{echte} Segmentquadratur (\file{mortarcontact.py:\allowbreak{}650--673}), sodass die Biorthogonalität
auf dem tatsächlichen Integrationsgebiet erfüllt ist (konsistente Randbehandlung nach
\citet{cichosz}, Abschnitt~\ref{sec:subcells}). Der Zusatzfaktor $\bm T_e$ ist die
Basistransformation für quadratische Elemente (Abschnitt~\ref{sec:basistrafo}); für lineare Elemente
ist $\bm T_e=\bm I$.

% =====================================================================
\section{Duale Basis für quadratische Elemente: die Basistransformation $\bm T_e$}
\label{sec:basistrafo}

Für quadratische Kontaktelemente ist die direkte duale Basiskonstruktion aus
Abschnitt~\ref{sec:dual} nicht unmittelbar brauchbar. Dieser Abschnitt begründet dies, leitet die
Abhilfe -- eine Basistransformation -- her und diskutiert die Wahl des Transformationsparameters
$\alpha$ einschließlich der begründeten Abweichung von der Literaturempfehlung
(vgl.\ Präsentation, Folie~9).

\subsection{Verlust der Integral-Positivität bei quadratischen Formfunktionen}
An die Multiplikator-Formfunktionen wird die \emph{Integral-Positivität} gestellt
(\citealp[Gl.~(4.2)]{popp2012}):
\begin{equation}
\int_{\text{ele}} \Phi_j\,\mathrm d\gamma \;>\; 0 .
\label{eq:intpos}
\end{equation}
Der Grund ist direkt algorithmisch: Ist \eqref{eq:intpos} verletzt, kann ein physikalisch offener
Spalt einen \emph{negativen} gewichteten Spalt ergeben (und umgekehrt); die Active-Set-Strategie
behandelt dann aktive Knoten als inaktiv und umgekehrt und konvergiert nicht
(\citealp{popp2012}: \emph{``the numerical algorithm will generate nonphysical gaps and penetrations \ldots\
a nonconverging active set strategy''}).

Genau das tritt bei quadratischen Formfunktionen auf. Das \emph{Eck}knoten-Integral ist nicht positiv
(\citealp[Gl.~(4.3)/(4.4)]{popp2012}):
\begin{equation}
\int_{\text{ele}} N_1^{\text{quad8}}\,\mathrm d\gamma < 0,
\qquad
\int_{\text{ele}} N_1^{\text{tri6}}\,\mathrm d\gamma = 0 .
\end{equation}
Ein anschauliches, numerisch greifbares Beispiel liefert \citet[Gl.~(29)]{puso2008}:
zwei ebene, \emph{nicht} durchdrungene Serendipity-Flächen im Abstand $l$ erzeugen am Eckknoten den
gewichteten Spalt $g_A = A\,l/3 > 0$ -- also eine \emph{vorgetäuschte} Durchdringung, allein wegen
der Nichtpositivität der quadratischen Funktion. (Ein wörtlicher Wert ``$-1/12$'' als negatives
\emph{Gewicht} kommt in der Literatur nicht vor; der im Code-Docstring genannte Zahlenwert bezieht
sich auf das negative Eckknoten-Integral im Sinne von \citet[Gl.~(4.3)]{popp2012}.)

\subsection{Basistransformation der Formfunktionen}
Statt der originalen Formfunktionen wird eine transformierte Basis verwendet
(\citealp[Gl.~(4.5)]{popp2012}):
\begin{equation}
\tilde{\bm N} = \bm T_e\,\bm N,
\end{equation}
deren $\bm T_e$ jede \emph{Kantenmittelknoten}-Zeile mit dem Anteil $\alpha$ auf die beiden
benachbarten Eckknoten umverteilt und $1-2\alpha$ auf sich behält; die Eckknotenzeilen bleiben die
Identität. Für CONQUAD8 lautet die Matrix (mit $\alpha=\tfrac13$) explizit
\begin{equation}
\bm T_e=\begin{bmatrix}
1&0&0&0&\tfrac13&0&0&\tfrac13\\
0&1&0&0&\tfrac13&\tfrac13&0&0\\
0&0&1&0&0&\tfrac13&\tfrac13&0\\
0&0&0&1&0&0&\tfrac13&\tfrac13\\
0&0&0&0&\tfrac13&0&0&0\\
0&0&0&0&0&\tfrac13&0&0\\
0&0&0&0&0&0&\tfrac13&0\\
0&0&0&0&0&0&0&\tfrac13
\end{bmatrix}
\qquad (\text{\file{element.py:\allowbreak{}424--468}}).
\end{equation}
Die Transformation hat drei nachweisbare Eigenschaften (\citealp[S.~B431]{popp2012}): sie ist
\emph{symmetrisch} (jeder Kantenknoten trägt gleich viel zu seinen beiden Eck-Nachbarn bei), sie
erhält die \emph{Partition of Unity} ($\sum_j \tilde N_j = 1$), und sie ist durch
$\alpha < \tfrac12$ nach oben beschränkt, damit auch die Kantenknoten-Integrale positiv bleiben. Aus
der transformierten Basis werden dann -- wie bei linearen Elementen über \eqref{eq:Ae} -- die dualen
Funktionen gebildet.

Die Transformation wird nur auf die \emph{serendipity}-Typen mit nichtpositivem Eckintegral
angewandt (CONQUAD8, CONTRI6; CONLINE3). Für das voll-Lagrange'sche CONQUAD9 (aus hex27) entfällt
sie: dessen Eckknoten-Integrale sind bereits strikt positiv ($\tfrac19$ auf dem Referenzquadrat),
und das Kanten-zu-Ecken-Schema erfasst den zusätzlichen Zentrumsknoten ohnehin nicht -- für CONQUAD9
wird daher wie für lineare Typen die Identität verwendet (\file{element.py:\allowbreak{}454--462}).

\subsection{Wahl des Transformationsparameters $\alpha$}
\label{sec:alphachoice}
Die \emph{Struktur} der Transformation (symmetrische Kanten-zu-Ecken-Umverteilung, Erhaltung der
Partition of Unity, obere Schranke $\alpha<\tfrac12$) geht auf \citet{lamichhane2005}
zurück, die die duale Basis für Elemente höherer Ordnung durch eine
\emph{elementweise Umverteilung} von Kanten-Mittelknoten auf die Eckknoten konstruieren
(\citealp[Gl.~(3.10)/(3.12)]{lamichhane2005}); \citet{popp2012} verallgemeinern dies auf die
Finite-Deformation. Der \emph{Zahlenwert} von $\alpha$ ist damit noch nicht festgelegt; die Literatur
bietet mehrere begründete Kandidaten:
\begin{itemize}\itemsep2pt
  \item \emph{Integral-Positivität} \eqref{eq:intpos} verlangt im unverzerrten Fall $\alpha>0$
        (tri6) bzw.\ $\alpha>\tfrac18$ (quad8).
  \item Die zusätzliche exakte Reproduktion linearer Polynome ($P_1$) im dualen Raum ergibt
        (unverzerrt) $\alpha=\tfrac{1}{12}$ (tri6) bzw.\ $\alpha=\tfrac15$ (quad8). Genau diese Werte
        leiten \citet{lamichhane2005} aus Biorthogonalität und $P_1$-Reproduktion her: den Faktor
        $\tfrac{1}{12}$ für quadratische Tetraeder (Gl.~(3.10)) und $\tfrac15$ für die
        Serendipity-Hexaeder, deren Facette gerade das quad8-Element ist (Gl.~(3.12)).
  \item \citet{popp2012} empfehlen für quadratische \emph{Flächen}elemente
        $\alpha=\tfrac15$, das nach ihrer Analyse die Integral-Positivität für alle praktisch
        auftretenden Elementverzerrungen sichert (\citealp[S.~B431]{popp2012}).
  \item \citet{farah} wählt für tet10-\emph{Volumen}elemente $\alpha_q=\tfrac13$
        (\citealp[Gl.~(6.20), S.~164]{farah}), begründet durch die Schranke $\alpha_q<\tfrac12$,
        Integral-Positivität und numerische Erfahrung.
\end{itemize}
Die vorliegende Implementierung verwendet $\alpha=\tfrac13$ (\file{element.py:\allowbreak{}451}) und weicht damit
bewusst von Popps flächenspezifischer Empfehlung $\tfrac15$ ab. Diese Wahl ist nicht kosmetisch,
sondern durch eine Eigenschaft der \emph{segmentbasierten} Auswertung begründet, die im Folgenden
dargelegt wird.

\paragraph{Integral-Positivität genügt bei partieller Überdeckung nicht.}
Popps Garantie bezieht sich auf die Positivität des \emph{Voll-Element}-Integrals
$\int_{\text{ele}}\tilde N_j\,\mathrm d\gamma$. In der segmentbasierten Formulierung
(Abschnitt~\ref{sec:subcells}) werden die dualen Koeffizienten und die knotenweisen Gewichte
$\tilde D_{jj}$ jedoch über die \emph{tatsächliche Überlappungsfläche} $\gamma_c^{\text{ovl}}\subseteq
\text{ele}$ integriert (konsistente Randbehandlung nach \citet{cichosz}). Ist ein Slave-Element
nur \emph{teilweise} in Kontakt, so ist das relevante Gewicht
$\int_{\gamma_c^{\text{ovl}}}\tilde N_j\,\mathrm d\gamma$; dessen Positivität für \emph{beliebige}
Teilgebiete ist genau dann gesichert, wenn die transformierte Formfunktion nicht nur ein positives
Integral besitzt, sondern \emph{punktweise nichtnegativ} ist,
\begin{equation}
\tilde N_j(\boldsymbol\xi) \;\ge\; 0 \quad \forall\,\boldsymbol\xi\in\text{ele},
\label{eq:pointwisepos}
\end{equation}
eine strikt stärkere Forderung als \eqref{eq:intpos}. Für die serendipity-Eckfunktion existiert eine
Schwelle $\alpha^\ast$, oberhalb derer \eqref{eq:pointwisepos} gilt; sie liegt zwischen den beiden
Literaturwerten $\tfrac15$ und $\tfrac13$.

\paragraph{Numerischer Befund.}
Dies ist im Verifikationslauf direkt belegt (Abschnitt~\ref{sec:verif}, Tests~05 und~08):
\begin{itemize}\itemsep2pt
  \item \textbf{Test~05 (partielle Überdeckung, CONQUAD8):} Bei um $0.3$ verschobenem Master
        ($70\%$ Überdeckung) sind mit $\alpha=\tfrac13$ alle knotenweisen Gewichte strikt positiv;
        mit $\alpha=\tfrac15$ wird ein Eckknoten-Gewicht \emph{negativ} ($\approx-0.0049$) --
        $\tilde N_{\text{Ecke}}$ ist dort also nicht punktweise nichtnegativ.
  \item \textbf{Test~08 (verzerrter hex20-Patch-Test):} Der übertragene Kontaktdruck weicht mit
        $\alpha=\tfrac13$ punktuell um $14.3\%$ vom Sollwert ab (innerhalb der $20\%$-Toleranz für
        verzerrte Netze), mit $\alpha=\tfrac15$ dagegen um $21.8\%$ (außerhalb der Toleranz). Die
        übertragene \emph{Gesamtkraft} ist in beiden Fällen exakt.
\end{itemize}
Der Wert $\alpha=\tfrac13$ liegt mit größerem Abstand zur unteren Positivitätsschranke und macht die
transformierten Eckfunktionen -- im untersuchten Bereich -- punktweise nichtnegativ, sodass die
Gewichte auch bei Teilkontakt und Netzverzerrung positiv und die diskrete Signorini-Bedingung robust
bleiben. Der Preis ist der Verzicht auf die exakte Reproduktion linearer Polynome, die $\alpha=\tfrac15$
auf unverzerrten Elementen böte; für die Konsistenzordnung des Kontakts ist dies unkritisch, und die
unverzerrten Patch-Tests werden mit beiden Werten exakt bestanden.

\begin{quote}\small
\textbf{Bemerkung.} Da im Zielanwendungsfall (Ausziehversuch eines Ankers im Betonblock)
verzerrte Netze und über den Lastverlauf wechselnder Teilkontakt auftreten, ist die durch
$\alpha=\tfrac13$ gesicherte punktweise Nichtnegativität hier die ausschlaggebende Eigenschaft --
sie hat Vorrang vor der exakten Linearreproduktion von $\alpha=\tfrac15$. Dies deckt sich mit Farahs
Begründung ``numerische Erfahrung'' (\citealp[Gl.~(6.20)]{farah}).
\end{quote}

\subsection{Auswirkung auf die Struktur von $\bm D$}
Mit der Transformation ist die globale Matrix $\bm D$ \emph{nicht} mehr diagonal.
\citet[Gl.~(4.12)]{popp2012} zeigen aber, dass sie sich exakt in zwei trivial invertierbare Faktoren spaltet:
\begin{equation}
\bm D = \tilde{\bm D}\,\bm T^{-1},
\qquad \tilde D_{jk} = \int \Phi_j \tilde N_k\,\mathrm d\gamma \ \text{diagonal},
\qquad \bm D^{-1} = \bm T\,\tilde{\bm D}^{-1}.
\end{equation}
Alle Vorteile der dualen Basis bleiben also erhalten, \emph{ohne} $\bm D$ zu ``lumpen'' (diagonal zu
approximieren). Ein naives Lumping wäre hier falsch: es verletzte die Zeilensummen-Identität
\eqref{eq:rowsum} bzw.\ die Konsistenzbedingung $\tilde D_{kk}=\sum_j M_{kj}$ und damit den
Patch-Test (dieser Konsistenz-Zusammenhang ist bei \citet[Abschn.~4.1]{cichosz} belegt).
Deshalb wird $\bm D$ im Code voll (nicht-gelumpt) assembliert (\file{mortarcontact.py:\allowbreak{}740--749}).

% =====================================================================
\section{Quadratische Facetten: lineare Sub-Cells}
\label{sec:subcells}

\subsection{Notwendigkeit der linearen Zerlegung}
Der Clipping-Algorithmus (Abschnitt~\ref{sec:geometrie}) schneidet Polygone in einer \emph{ebenen}
Hilfsfläche und ist nur für \emph{gerade} Kanten korrekt. Eine quadratische Facette (CONQUAD8/9,
CONTRI6) ist im Allgemeinen gekrümmt und hat quadratisch interpolierte Kanten.
\citet{puso2008} lösen das, indem die gekrümmte Facette in \emph{linear interpolierte} Sub-Segmente
zerlegt wird, auf die der Schnitt-Algorithmus unverändert anwendbar ist:
\begin{quote}
\small ``our approach to geometrically approximate these quadratic surfaces will be to subdivide them
into linearly interpolated contact sub-segments \ldots\ The algorithms described \ldots\ for
intersection of linearly defined element surfaces can then be applied almost without modification. In
so doing, we of course sacrifice any type of geometrically exact description of the quadratic
contacting surfaces in defining the integration algorithm.'' \hfill \citep[S.~560--561]{puso2008}
\end{quote}

\subsection{Die konkrete Zerlegung}
Die Zerlegung entspricht exakt \citet[Abb.~3]{puso2008}: das Serendipity-Element (quad8) wird in
\emph{vier lineare Eck-Dreiecke und ein zentrales Quadrilateral} zerlegt (Abb.~3(e)/(f)), quad9 in
vier Quads, tri6 in vier lineare Dreiecke (Abb.~3(c)/(d)). Im Code steht dies als \code{SUB\_CELL\_MAP}
(\file{mortarcontact.py:\allowbreak{}156--183}):
\begin{center}
\small
\begin{tabular}{ll}
\toprule
Elementtyp & Sub-Cells (lokale Knotenindizes) \\
\midrule
CONQUAD8 & $[0,4,7],\;[4,1,5],\;[5,2,6],\;[7,6,3]$ und Mittel-Quad $[4,5,6,7]$ \\
CONQUAD9 & $[0,4,8,7],\;[4,1,5,8],\;[8,5,2,6],\;[7,8,6,3]$ \\
CONTRI6  & $[0,3,5],\;[3,4,5],\;[3,1,4],\;[5,4,2]$ \\
CONLINE3 & $[0,2],\;[2,1]$ (2D-Analogon) \\
\bottomrule
\end{tabular}
\end{center}
Die Knoten $4$--$7$ (bzw.\ $3$--$5$) sind die Kantenmittelknoten; die Sub-Cells werden über sie
aufgespannt.

\subsection{Wie die Krümmung dennoch eingeht}
Wesentlich: die Zerlegung betrifft \emph{nur} das Integrationsgebiet (die Clip-Geometrie), \emph{nicht}
die Feldinterpolation. Über affine Abbildungen zwischen Sub-Segment- und Eltern-Element-Parameterraum
(\citealp[Gl.~(31)/(32)]{puso2008}) werden an den Gauß-Punkten weiterhin die \emph{vollen}
quadratischen Formfunktionen ausgewertet:
\begin{quote}
\small ``it is possible to evaluate higher-order shape function products in the integrand of the
mortar matrices and the weighted gap without any algorithmic changes. This approach only affects the
integration domain itself, which is less accurate in terms of geometry.''
\hfill (\citealp[Anh.~A.1.4, S.~215]{farah}; Verweis auf \citealt{puso2008})
\end{quote}
Die Krümmung geht also ausschließlich über die Formfunktions-Auswertung an den (per Newton
rückabgebildeten) Gauß-Punkten ein, nie über die geraden Clip-Kanten. Dieses Vorgehen ist inhaltlich
identisch zu MOOSEs \code{AutomaticMortarGeneration} (``Linearize secondary face elements''). Damit
wird der hex20-Patch-Test exakt bestanden (Abschnitt~\ref{sec:verif}).

% =====================================================================
\section{Knotennormalen und eingefrorene Geometrie}
\label{sec:normalen}

\subsection{Flächengewichtete Knotennormale}
Der gewichtete Spalt \eqref{eq:wgap_pos} und die Projektion benötigen eine eindeutige Normale pro
Slave-Knoten. Sie wird als \emph{flächengewichtetes Mittel} der angrenzenden Facettennormalen gebildet
und normiert (\file{mortarcontact.py:\allowbreak{}402--458}):
\begin{equation}
\bm n_I = \frac{\sum_{f\ni I}\bm n_f}{\big\|\sum_{f\ni I}\bm n_f\big\|},
\qquad
\bm n_f = \tfrac12(\bm v_1-\bm v_0)\times(\bm v_2-\bm v_0)
\end{equation}
(im 2D: $\bm n=[t_y,-t_x]$ aus der Kantentangente). Diese gemittelte Normale garantiert ein
stetiges Normalenfeld über Facettengrenzen hinweg, was für die Konsistenz des Kontakts an
Elementkanten nötig ist.

\subsection{Eingefrorene Geometrie (staggered update)}
Normalen $\bm n_I$ und Koppelmatrizen $\bm D,\bm M$ werden \emph{einmal} zu Increment-Beginn in der
aktuellen (deformierten) Konfiguration berechnet und dann für alle Newton-Iterationen des Increments
festgehalten (\file{mortarcontact.py:\allowbreak{}738}). Dadurch ist die assemblierte Steifigkeit die
\emph{exakte} Jacobi-Matrix des Residuums bei fixierter Geometrie -- Voraussetzung für die
Newton-Konvergenz (Abschnitt~\ref{sec:solve}). Der Preis ist eine gestaffelte (staggered)
Behandlung der Geometrienichtlinearität, die durch hinreichend kleine Lastschritte kontrolliert
wird.

\medskip
\noindent\textit{Bekannte Einschränkung.} An \emph{scharfen} Kanten (z.\,B.\ Übergang
Ankerschaft$\to$Ankerkopf im hex20-Netz) ist die gemittelte Knotennormale mehrdeutig
($\sim45^\circ$), was die Projektion dort verschlechtern kann. Für die hier betrachteten glatten
Kontaktflächen ist das unkritisch.

% =====================================================================
\section{Signorini als semismooth NCP}
\label{sec:ncp}

Die diskreten Kontaktbedingungen \eqref{eq:kkt} sind Ungleichungen mit Fallunterscheidung
(Kontakt / offen). Dieser Abschnitt fasst sie in \emph{eine} nichtglatte Gleichung -- die
NCP-Funktion (Folie~11) -- und begründet ihre Lösung per Active-Set/semismooth Newton.

\subsection{Von den KKT-Bedingungen zur NCP-Funktion}
Die knotenweisen Hertz--Signorini--Moreau-Bedingungen
$\tilde g_{j}\ge0,\; \lambda_{n,j}\ge0,\; \lambda_{n,j}\tilde g_{j}=0$
(\citealp[Gl.~(3.41)]{farah}; \citealp[Gl.~(38)]{gitterle}) lassen sich mit einem beliebigen
$c_n>0$ äquivalent als \emph{eine} nichtglatte Komplementaritätsfunktion (NCP) schreiben. Die Urform
stammt aus \citet[Gl.~(3.9)]{hueber2005}; in der Kontakt-Notation lautet sie
(\citealp[Gl.~(55)]{gitterle}; \citealp[Gl.~(3.58)]{farah})
\begin{equation}
C_{n,j}\big(\lambda_{n,j},\bm d\big) \;=\; \lambda_{n,j} - \max\!\big(0,\; \lambda_{n,j} - c_n\,\tilde g_{j}\big) \;=\; 0,
\qquad c_n>0 .
\label{eq:ncp}
\end{equation}
\textbf{Äquivalenz zur KKT-Form.} Mit der elementaren Identität
$a - \max(0, a-b) = \min(a,b)$ ist \eqref{eq:ncp} genau die \emph{min}-NCP-Funktion
\begin{equation}
C_{n,j}=0 \;\Longleftrightarrow\; \min\!\big(\lambda_{n,j},\; c_n\,\tilde g_{j}\big)=0
\;\Longleftrightarrow\; \big[\,\lambda_{n,j}\ge0,\;\ \tilde g_{j}\ge0,\;\ \lambda_{n,j}\,\tilde g_{j}=0\,\big],
\end{equation}
also exakt die drei KKT-Bedingungen -- und zwar \emph{unabhängig} vom Wert $c_n$
(\citealp[Gl.~(56)]{gitterle}: erfüllt (38) \emph{``independently of the choice of the parameter $c_n$''}). (Die
min/max-Identität ist Standard; eine Fischer--Burmeister-Formulierung wird in den Quellen nicht
verwendet und hier nicht benötigt.)

\medskip
\noindent\textbf{Wer und warum.} Die Normal-NCP wurde für Kleindeformation von
\citet{hueber2005} eingeführt und von \citet{poppgeewall2009} auf
finite Deformation übertragen -- dort erscheint sie im reibungsfreien Fall in genau dieser Form
(\citealp[Gl.~(56)]{poppgeewall2009}); \citet{gitterle} fassen dies in Gl.~(55) zusammen. Die
mathematische Wurzel der später genutzten Äquivalenz ``Active-Set $=$ semismooth Newton'' ist
\citet{hik2002} (Satz~3.1: lokal superlineare, bei starker
Semismoothness sogar $q$-quadratische Konvergenz).

\subsection{Der Aktivierungsindikator und die Codeform}
Aktiv/inaktiv wird über das Vorzeichen des Arguments der $\max$-Funktion entschieden
(\citealt[Gl.~(3.13)]{hueber2005} bzw.\ \citealt[Gl.~(70)]{gitterle} -- letztere äquivalent als komplementäre \emph{Inaktiv}-Menge mit $\le 0$ geschrieben):
\begin{equation}
\text{Knoten } j \text{ aktiv} \iff s_{n,j} := \lambda_{n,j} - c_n\,\tilde g_{j} > 0 .
\label{eq:indicator}
\end{equation}
Im Code steht diese Bedingung mit einer \emph{zusätzlichen Normierung} mit $\tilde D_{jj}^{-1}$
(\file{mortarcontact.py:\allowbreak{}797--839}):
\begin{equation}
s_{n,j} = p_{n,j} - c_n\,\tilde D_{jj}^{-1}\,\tilde g_{j},
\qquad \text{aktiv} \iff s_{n,j}>0,
\label{eq:indicator_code}
\end{equation}
mit dem physikalischen Druck $p_{n,j}=-\lambda_j\,\mathrm{sgn}(D)$. \emph{Bemerkung:} der
Faktor $\tilde D_{jj}^{-1}$ ist \emph{nicht} Teil der Originalgleichung~(55)/(3.58) -- dort steht
$c_n\tilde g_j$ mit dem \emph{gewichteten} Spalt direkt. Da $\tilde g_j$ ein volumetrisches,
über $\tilde D_{jj}$ skaliertes Maß ist (Abschnitt~\ref{sec:koppelmatrizen}), macht die Division
durch $\tilde D_{jj}$ den Term $c_n\tilde D_{jj}^{-1}\tilde g_j$ \emph{druckartig} und damit
dimensionell vergleichbar mit $p_{n,j}$. Es ist eine bewusste, dimensionell konsistente
\emph{Normierungsentscheidung} der Implementierung (im Einklang mit der $D^{-1}$-Skalierung der
dualen Mortar-Methode), kein Wortzitat aus Gitterle. Die Struktur der NCP und der
Aktivierungsindikator \eqref{eq:indicator} sind dadurch unberührt.

\subsection{Der Parameter $c_n$}
$c_n>0$ ist ein rein \emph{algorithmischer} Parameter: bei Konvergenz ist entweder $\tilde g_j=0$
(aktiv) oder $\lambda_{n,j}=0$ (inaktiv), sodass $c_n$ aus $\min(\lambda_{n,j},c_n\tilde g_j)$
herausfällt -- er beeinflusst \emph{die Lösung nicht}, nur das Konvergenzverhalten
(\citealp[S.~565]{gitterle}: \emph{``purely algorithmic parameters with no influence on the accuracy of
results''}; \citealp[S.~38]{farah}). Für die Größenordnung wird
\begin{equation}
c_n \sim \mathcal O(E)
\end{equation}
empfohlen (\citealp[S.~1367]{poppgeewall2009}: \emph{``$c_{\mathrm n}$ is suggested to be at the order
of Young's modulus $E$ of the contacting bodies''}; \citealp[S.~38]{farah}; \citealp[S.~565]{gitterle}:
Parameter ``in the range of material parameters''). Der Grund ist eine Skalenbalance: $\lambda_{n,j}$
ist druckartig ($\sim E\cdot$Dehnung), $\tilde g_j$ längen-/volumenartig; $c_n\sim E$ (bzw.\ nach der
Normierung \eqref{eq:indicator_code} $c_n/\tilde D_{jj}\sim E$) bringt beide auf dieselbe Skala. Zu
\emph{großes} $c_n$ destabilisiert dagegen die Active-Set-Iteration: \citet[Tab.~I, S.~565]{gitterle}
zeigt für $c_n\in\{10^5,10^6\}$ mehrfache Aktiv-Mengen-Wechsel (\emph{Chattering}, dort mit $^*$
markiert) bis hin zur Nichtkonvergenz, während ein breites mittleres Band robust in wenigen Schritten
konvergiert. Im Code ist $c_n$ mit $\tilde D_{jj}$ normiert und in der Größenordnung von $E$ gewählt
(\file{mortarcontact.py:\allowbreak{}92--108}), mit demselben Chattering-Argument im Kommentar.

\subsection{PDASS als semismooth Newton}
Die $\max$-Funktion in \eqref{eq:ncp} ist \emph{semismooth}; ihre verallgemeinerte Ableitung ist
$\partial_x\max(a,x)=\{0 \text{ falls } x\le a,\ 1 \text{ falls } x>a\}$ (\citealp[Gl.~(64)]{gitterle}).
Damit ist die Primal-Dual Active Set Strategy (PDASS) ein \emph{semismooth Newton-Verfahren}: das
Active Set wird in jeder Iteration aus \eqref{eq:indicator} neu bestimmt, das resultierende lineare
System (bei fixiertem Set glatt) gelöst, und iteriert, bis sich das Set nicht mehr ändert
(Stopp-Kriterium, \citealp[S.~3153]{hueber2005}; \citealp[Abschn.~3.5.2]{farah}, ``re-interpreted as a
semi-smooth Newton method''). Die abstrakte Grundlage dieser Interpretation und Konvergenz liefert
\citet{hik2002} (Satz~3.1); den ersten \emph{konsistent} linearisierten
Newton-Ansatz für den dualen Mortar-Kontakt mit PDASS setzte \citet[Gl.~(65)/(67)]{poppgeewall2009} um. Die beobachtete Konvergenz ist charakteristisch: solange das
Set noch wechselt, sinkt das Residuum nur langsam; \emph{sobald das Set fixiert ist, konvergiert das
Verfahren quadratisch} -- als Folge der konsistenten Linearisierung (\citealp[S.~560]{gitterle};
\citealp[Tab.~4.1]{farah}, mit Residuenabfall $\sim10^{-3}\!\to\!10^{-8}\!\to\!10^{-11}$ nach Fixierung). Die
Implementierung bestätigt das Set-Stabilitäts-Kriterium und sichert es zusätzlich per Anti-Cycling
(Bland-Regel) gegen Zustandswiederholung ab (\file{mortarcontact.py:\allowbreak{}841--867}; siehe
Abschnitt~\ref{sec:solve}).

% =====================================================================
\section{Lösungsstrategie: Newton-Ablauf und Active-Set}
\label{sec:solve}

\subsection{Sattelpunkt-Struktur}
Auf Branch \code{0\_mortarcontact} sind die Lagrange-Multiplikatoren \emph{explizite} zusätzliche
Unbekannte (ein skalarer Normalmultiplikator pro Slave-Knoten,
\file{mortarcontact.py:\allowbreak{}348--351}). Das globale System ist damit ein \emph{Sattelpunkt}-System
\begin{equation}
\begin{bmatrix} \bm K & \bm B^\top \\ \bm B & \bm 0 \end{bmatrix}
\begin{bmatrix} \Delta\bm d \\ \Delta\bm\lambda \end{bmatrix}
= -\begin{bmatrix} \bm r_d \\ \bm r_\lambda \end{bmatrix},
\qquad \bm B = [\,\bm D\ \ -\bm M\ \ \bm 0\,],
\end{equation}
mit den Kontaktkräften aus \eqref{eq:DM}. Eine statische Kondensation der Multiplikatoren
($\bm P=\bm D^{-1}\bm M$) ist auf diesem Branch bewusst nicht aktiviert.

\subsection{Ein Newton pro Increment mit eingefrorener Geometrie}
Material-, Geometrie- und Kontaktnichtlinearität werden in \emph{einem} Newton-Verfahren pro
Lastschritt gelöst. Wie in Abschnitt~\ref{sec:normalen} beschrieben, werden Normalen und $\bm D,\bm M$
zu Increment-Beginn eingefroren; die assemblierte Tangente ist dann die exakte Jacobi-Matrix des
Residuums, was die (nach Fixierung des Active Set) quadratische Konvergenz ermöglicht. Die
konsistente Tangente ist per Finite-Differenzen zu $\sim10^{-10}$ verifiziert
(Abschnitt~\ref{sec:verif}).

\subsection{Active-Set-Iteration und Anti-Cycling}
Das Active Set wird in \emph{jeder} Newton-Iteration aus dem Indikator \eqref{eq:indicator_code} neu
bestimmt (\file{mortarcontact.py:\allowbreak{}797--839}). Pro Knoten gibt es zwei Zweige:
\begin{itemize}\itemsep2pt
  \item \textbf{aktiv} ($s_{n,j}>0$): Constraint-Residuum $\text{PExt}[\lambda]\!-\!\!=\!\tilde g_j$;
        Slave-Kräfte $+\lambda_j\bm D[j,k]\bm n_j$, Master-Kräfte $-\lambda_j\bm M[j,l]\bm n_j$;
  \item \textbf{inaktiv} ($s_{n,j}\le0$): $\lambda_j$ wird zu null gezwungen
        ($\text{PExt}[\lambda]\!-\!\!=\!\lambda_j$, $K[\lambda,\lambda]\!+\!\!=\!1$).
\end{itemize}
Die äußere Iteration ist konvergiert, sobald sich das Set nicht mehr ändert (PDASS-Kriterium,
\citealp{hueber2005}). Gegen \emph{Zyklen} (abwechselnde Active-Set-Zustände) sichert eine
Anti-Cycling-Regel im Sinne von Bland: bei erneutem Besuch eines bereits gesehenen Zustands wird das
Set eingefroren (\file{mortarcontact.py:\allowbreak{}841--867}); eine harte Obergrenze von 20 Iterationen dient als
Sicherheitsnetz.

% =====================================================================
\section{Verifikation}
\label{sec:verif}

Die Implementierung ist durch eine Testsuite unter \file{testfiles/mortar\_tests/} abgesichert; die
folgenden Bausteine belegen die Korrektheit der obigen Herleitungen.

\begin{center}
\small
\begin{tabularx}{\linewidth}{lX}
\toprule
Eigenschaft & Prüfung \\
\midrule
\textbf{Polygon-Clipping} & Schnitt-/Triangulierungsgeometrie gegen analytische Überlappungen (\file{02\_polygon\_clipping}). \\
\textbf{Koppelmatrizen} & $\bm D,\bm M$ und Zeilensummen-Identität \eqref{eq:rowsum} numerisch (\file{03\_coupling\_matrices}). \\
\textbf{Biorthogonalität} & \eqref{eq:biorth} numerisch verifiziert (\file{04\_dual\_biorthogonality}). \\
\textbf{Quadratische Segmentierung} & Sub-Cell-Zerlegung + Segmentquadratur für CONQUAD8/9, CONTRI6 (\file{05\_quadratic\_segmentation}). \\
\textbf{Active-Set / PDASS} & Aktivierung bei Durchdringung, Freeze/Anti-Cycling (\file{06\_active\_set\_pdass}). \\
\textbf{Patch-Test 2D} & konstanter Druck exakt übertragen, CPE4/CPE8, konform + nichtkonform (\file{07\_patch\_test\_2d}). \\
\textbf{Patch-Test hex20} & 3D-Patch-Test mit CONQUAD8 exakt bestanden (\file{08\_patch\_test\_hex20}). \\
\bottomrule
\end{tabularx}
\end{center}

\noindent Ergänzend gilt:
\begin{itemize}\itemsep2pt
  \item \textbf{Kontakt-Patch-Test} (konstanter Druck exakt übertragen) in 2D (CPE4/CPE8) und 3D
        (hex8/hex20), konform \emph{und} nichtkonform. Das bestätigt zugleich die
        Zeilensummen-Identität \eqref{eq:rowsum} und die korrekte Behandlung quadratischer Elemente
        (Basistransformation, Sub-Cells).
  \item \textbf{Konsistente Tangente} per Finite-Differenzen bestätigt ($\sim10^{-10}$) -- die
        Voraussetzung für die quadratische Newton-Konvergenz nach Fixierung des Active Set.
  \item \textbf{Zwei-Block-Referenztests} (\file{Control\_Tests/README.md}): verschiebungs- und
        lastgesteuerte Blöcke, konform und nichtkonform, mit Abgleich von Verschiebungsfeld,
        Spannungen und Kontaktbedingungen.
\end{itemize}
Der reibungsfreie Mortar-Kontakt rechnet damit stabil -- Increment für Increment bei vollem
Lastschritt.

% =====================================================================
\section{Forschung: Warum ist \code{hex20} im Ausziehversuch steifer als \code{hex8}?}
\label{sec:forschung}

Der Ausziehversuch eines Kopfbolzendübels aus einem Betonquader (\file{testfiles/mortar\_tests/\allowbreak POT\_Dejori})
wird mit demselben Modell einmal mit \emph{linearen} Elementen (\code{hex8}: Beton \code{GC3D8},
Kontakt \code{CONQUAD4}) und einmal mit \emph{quadratischen} Elementen (\code{hex20}: Beton
\code{GC3D20R}, Kontakt \code{CONQUAD8}) gerechnet. Beide Netze sind bis auf die zusätzlichen
Kantenmittelknoten identisch (gleiche Elementzahl, dieselben sechs Kontaktflächen), das Material
(gradientenerweiterte Schädigungsplastizität, \code{GCDP}) ist dasselbe, und beide Rechnungen sind
verschiebungsgesteuert. Dennoch liefern sie deutlich verschiedene Kraft-Weg-Kurven. Da die
Verifikation (Abschnitt~\ref{sec:verif}) die Korrektheit der Implementierung für lineare \emph{und}
quadratische Elemente belegt, stellt sich die Frage: Wie kann dieser Unterschied zustande kommen?

\subsection{Beobachtung}
Der Unterschied ist in erster Linie eine \emph{Steifigkeitsdifferenz}, und sie ist bereits im
nahezu elastischen Bereich vollständig vorhanden -- die Sekantensteifigkeit ist über die ersten
$0{,}03$~mm praktisch konstant, also \emph{bevor} nennenswerte Schädigung eintritt und \emph{bevor}
das Active-Set aktiv wird.

\begin{center}
\small
\begin{tabular}{lcc}
\toprule
 & \code{hex8} (linear) & \code{hex20} (quadratisch) \\
\midrule
Sekantensteifigkeit ($u\le 0{,}03$~mm) & $\approx 117$~kN/mm & $\approx 197$~kN/mm \\
Peak-Last                              & $22{,}6$~kN         & $19{,}9$~kN \\
Weg beim Peak                          & $0{,}22$~mm         & $0{,}12$~mm \\
\bottomrule
\end{tabular}
\end{center}

\noindent Im Vergleich mit den drei Laborversuchen trifft \code{hex20} sowohl die Anfangssteifigkeit
als auch die Höhe des Peaks, während \code{hex8} zu weich ist und den Peak überschätzt. Es ist also
nicht so, dass \code{hex20} \emph{zu steif} wäre; vielmehr ist \code{hex8} \emph{zu weich}, und genau
das ist zu erklären. Weil die Differenz elastisch (vor Schädigung, vor Active-Set) auftritt, ist sie
weder eine Frage des Schädigungsmodells noch der Kontakt-Aktivierung.

\subsection{Der entscheidende Hinweis: die Steifigkeit bestimmt der Kontakt, nicht das Volumenelement}
Bei einem \emph{reinen} verschiebungsbasierten Finite-Elemente-Ansatz ist die diskrete Steifigkeit
eine obere Schranke der wahren Steifigkeit, und ein reicherer Ansatzraum senkt sie: quadratische
Elemente sind dort stets \emph{weicher} und genauer als lineare. Wäre also allein der Betonkörper
maßgeblich, müsste \code{hex20} weicher sein als \code{hex8} -- niemals steifer. Auch die konkrete
Integration der Volumenelemente zeigt in dieselbe Richtung: \code{GC3D8} ist voll integriert (und
damit eher zu steif, versteifungsanfällig), \code{GC3D20R} ist reduziert integriert (und damit eher
weicher). Beide Effekte auf der Volumenebene sagen \code{hex8}~$\ge$~\code{hex20} voraus; beobachtet
wird das Gegenteil.

Daraus folgt: Die zusätzliche Steifigkeit von \code{hex20} stammt aus der \emph{Kontakt}-Lasteinleitung,
nicht aus dem Volumenelement. Der Mortar-Kontakt ist eine \emph{gemischte} (Lagrange-)Formulierung, in
der die Monotonie ``reicherer Raum $=$ weicher'' gerade nicht gilt; und das grobe lineare Netz ist an
der Lasteinleitung erkennbar noch nicht konvergiert.

\subsection{Der physikalische Mechanismus: Lasteinleitung am Bolzenkopf}
Physikalisch ist der Ausziehversuch ein Lasteinleitungs- bzw.\ Stempelproblem: Der steife
Stahl\emph{kopf} presst sich gegen den Beton, und die Steifigkeit der Antwort hängt fast
ausschließlich davon ab, wie gut diese \emph{konzentrierte} Lasteinleitung am Kopf aufgelöst wird. Die
lineare Diskretisierung versagt dabei an zwei Stellen gleichzeitig:
\begin{itemize}\itemsep2pt
  \item \textbf{Randkonzentration des Kontaktdrucks.} Unter einem steifen Stempel steigt der
        Kontaktdruck zu den Rändern der Kontaktfläche hin stark an. Der lineare Mortar-Kontakt
        (diagonale $\bm D$, Multiplikatoren nur an Eckknoten) kann diesen Randanstieg nicht
        darstellen und verschmiert den Druck; die Grenzfläche gibt dort zu leicht nach -- die
        Antwort wird künstlich weich. Der quadratische Mortar-Kontakt (volle, nicht gelumpte
        $\bm D$, zusätzliche Kantenmittelknoten, transformierte duale Basis nach
        Abschnitt~\ref{sec:basistrafo}; \citealp{popp2012}) löst die Druckverteilung samt
        Randkonzentration deutlich besser auf.
  \item \textbf{Dehnungsgradient im Beton am Kopf.} Direkt am Kopf herrscht ein steiler
        Dehnungsgradient (hohe Pressung unter dem Kopf, rasch abklingend nach außen). Lineare
        Elemente verschmieren diesen Gradienten über das Element und mischen das hoch beanspruchte,
        steife Material am Kopf mit dem weicheren daneben zu einer zu nachgiebigen Mittelung;
        quadratische Elemente bilden den Gradienten ab.
\end{itemize}
Beides macht \code{hex8} zu weich und erklärt zugleich den überschätzten Peak: Weil die lineare
Diskretisierung die lokale Spannungsspitze am Kopf unterschätzt, erreicht der Beton seine Festigkeit
erst später; die Schädigung setzt verspätet ein, der Anker wird weiter gezogen, und die Kurve läuft
auf einen höheren, späteren Peak. Die quadratische Diskretisierung erfasst die lokale Spitze korrekt,
sodass die Schädigung bei der realen Last einsetzt und Peak-Höhe wie Peak-Lage zu den Versuchen passen.

\subsection{Abgleich mit der Verifikation}
Dies steht nicht im Widerspruch zu den bestandenen Patch- und Referenztests
(Abschnitt~\ref{sec:verif}). Diese prüfen gleichförmige Spannungs- bzw.\ Dehnungszustände an glatten
Grenzflächen unter Verschiebungssteuerung -- also gerade \emph{ohne} Konzentration. In diesem Regime
ist die Lösung elementordnungs\emph{unabhängig} und der Mortar-Kontakt exakt, für lineare wie für
quadratische Elemente. Der Ausziehversuch ist das gegenteilige Regime: Er ist
konzentrationsdominiert, und genau dann entscheidet die Elementordnung. Beide Diskretisierungen sind
korrekt; sie sind lediglich unterschiedlich gut aufgelöste Näherungen desselben, ordnungsempfindlichen
Problems.

\subsection{Schlussfolgerung}
\code{hex20} ist die genauere, nahezu konvergierte Lösung und trifft die Laborversuche; \code{hex8}
löst die konzentrierte Lasteinleitung am Bolzenkopf zu grob auf und ist deshalb künstlich weich, mit
überschätztem Peak. Der große Unterschied zwischen den beiden Kurven ist damit ein
Diskretisierungs- bzw.\ Elementordnungseffekt an der Lasteinleitung -- kein Implementierungsfehler.

% =====================================================================
\section{Zusammenfassung}
\label{sec:zusammenfassung}

Der auf Branch \code{0\_mortarcontact} implementierte reibungsfreie Mortar-Kontakt überträgt Kräfte
zwischen nichtkonformen Netzen über eine schwach (im integralen Mittel) durchgesetzte
Nichtdurchdringungsbedingung. Die tragenden Bausteine sind:
\begin{itemize}\itemsep2pt
  \item eine \textbf{segmentbasierte Integration} mit BVH-Nachbarsuche, Hilfsebenen-Projektion,
        Sutherland--Hodgman-Clipping und Grad-5-Quadratur (7-Punkt-Dreiecksregel nach \citealt{dunavant1985}),
        wobei der erhöhte Grad die nichtlineare Projektion auffängt (\citealp{farah2015});
  \item \textbf{duale Lagrange-Multiplikatoren} (\citealp{wohlmuth2000}), die $\bm D$ diagonalisieren und die
        Koppelmatrizen $\bm D,\bm M$ samt gewichtetem Spalt $\tilde g_j$ liefern, mit garantierter
        Zeilensummen-Identität (Impulserhaltung);
  \item eine \textbf{Basistransformation} für quadratische Elemente (\citealp{popp2012}; im Code der
        Farah-Wert $\alpha=\tfrac13$) zur Sicherung positiver dualer Gewichte, plus lineare
        \textbf{Sub-Cells} fürs Clipping (\citealp{puso2008});
  \item die \textbf{semiglatte NCP-Formulierung} der Signorini-Bedingung (\citealp[Gl.~(55)]{gitterle})
        gelöst per \textbf{PDASS}/semismooth Newton (\citealp{hueber2005}).
\end{itemize}
Die Formulierung besteht die Kontakt-Patch-Tests in 2D und 3D (linear und quadratisch, konform und
nichtkonform) und rechnet stabil bis in den entfestigenden Bereich des Betonmaterials. Reibung und
statische Kondensation sind nicht Teil dieses Branches und wurden hier bewusst ausgeklammert.

% =====================================================================
\section{Anwendungsanleitung: Kontaktspezifikation in der Input-Datei}
\label{sec:inputfile}

\subsection{Syntax-Struktur des Mortar-Constraint-Blocks}

Der Kontakt wird im Modell-Header unter \code{*constraint} definiert. Auf Branch
\code{0\_mortarcontact} kennt der Constraint vier Argumente (\file{mortarcontact.py:\allowbreak{}89--108}): die
beiden Pflicht-Oberflächen, das Feld und den NCP-Parameter \code{cn}:

\begin{lstlisting}[language=bash,caption={Mortar-Kontakt-Block (reibungsfrei) in der EdelweissFE Input-Datei}]
*constraint, type=mortarcontact, name=mein_kontakt_name
nonMortarSurface=slave_surface_name
mortarSurface=master_surface_name
field=displacement
cn=30000.0
\end{lstlisting}

Der Wert von \code{cn} sollte in der Größenordnung des Elastizitätsmoduls des weicheren
Kontaktpartners liegen (hier beispielhaft $\approx3\cdot10^4$ für Beton in MPa); die genaue
Bedeutung und Begründung stehen in Abschnitt~\ref{sec:ncp}.

\subsection{Schlüsselwörter}

\begin{table}[h]
\centering
\begin{tabularx}{\linewidth}{lX}
\toprule
\textbf{Schlüsselwort} & \textbf{Beschreibung} \\
\midrule
\code{type=mortarcontact} & Pflichtangabe. Aktiviert den segmentbasierten Mortar-Kontakt-Constraint (2D und 3D). \\
\code{name=\path{<string>}} & Eindeutiger Name des Kontaktpaars. \\
\code{nonMortarSurface=\path{<string>}} & Pflicht. Name der \textbf{Non-Mortar-Fläche (Slave)}. Trägt die Kontaktknoten, Normalen $\bm n_I$ und dualen Formfunktionen $\bm\Phi$. \\
\code{mortarSurface=\path{<string>}} & Pflicht. Name der \textbf{Mortar-Fläche (Master)}, auf die der Spalt gemessen wird. \\
\code{field=displacement} & Optional (Standard \code{displacement}). Verschiebungs-Vektorfeld, auf das der Kontakt wirkt. \\
\code{cn=\path{<float>}} & Optional (Standard $10^6$). Komplementaritätsparameter $c_n$ der Normal-NCP (Abschnitt~\ref{sec:ncp}). \emph{Sollte} pro Problem in der Größenordnung $\mathcal O(E)$ des weicheren Kontaktpartners gesetzt werden; der Default ist nur ein Rückfallwert. \\
\bottomrule
\end{tabularx}
\caption{Parameter des \code{mortarcontact}-Constraints auf Branch \code{0\_mortarcontact}.}
\end{table}

\subsection{Vorbereitung der Oberflächen und Kontaktelemente}

Der Mortar-Kontakt basiert auf expliziten Overlay-Kontaktelementen (\code{CONLINE2/3} in 2D bzw.\
\code{CONQUAD4/8/9}, \code{CONTRI3/6} in 3D). Beim Erzeugen aus Abaqus/Cubit-Netzen mit
\code{prepare\_mortar\_mesh.py} gilt:
\begin{itemize}
    \item \textbf{Normalenorientierung:} Elementnormalen müssen aus dem Kontinuumskörper nach außen
          weisen (\code{prepare\_mortar\_mesh.py} stellt dies automatisch sicher).
    \item \textbf{Slave/Master-Wahl:} Als Slave-Seite (\code{nonMortarSurface}) das feinere Netz
          bzw.\ das weichere Material wählen (Standardempfehlung der dualen Mortar-Methode).
\end{itemize}

% =====================================================================
\clearpage
\begin{thebibliography}{Hintermüller et al.(2002)}

\bibitem[Cichosz \& Bischoff(2011)]{cichosz} T.~Cichosz, M.~Bischoff: \emph{Consistent treatment of boundaries with mortar contact formulations using dual Lagrange multipliers}, Computer Methods in Applied Mechanics and Engineering 200 (2011), 1317--1332. (Lokale Kopie: \path{meetings/literature/CichoszBischoff_2011_ConsistentTreatmentOfBoundariesWithMortarContactFormulationsUsingDualLagrangeMultipliers.pdf})

\bibitem[Dunavant(1985)]{dunavant1985} D.~A.~Dunavant: \emph{High degree efficient symmetrical Gaussian quadrature rules for the triangle}, International Journal for Numerical Methods in Engineering 21 (1985), 1129--1148. (Lokale Kopie: \path{meetings/literature/Dunavant_1985_HighDegreeEfficientSymmetricalGaussianQuadratureRulesForTheTriangle.pdf})

\bibitem[Ericson(2004)]{ericson2004} C.~Ericson: \emph{Real-Time Collision Detection}, Morgan Kaufmann, 2004. (Lokale Kopie: \path{meetings/literature/Ericson_2004_RealTimeCollisionDetection.pdf})

\bibitem[Farah(2018)]{farah} P.~W.~Farah: \emph{Mortar Methods for Computational Contact Mechanics Including Wear and General Volume Coupled Problems}, Dissertation, TU M\"unchen, 2018. (Lokale Kopie: \path{meetings/literature/Farah_2018_MortarMethodsForComputationalContactMechanicsIncludingWearAndGeneralVolumeCoupledProblems.pdf})

\bibitem[Farah et al.(2015)]{farah2015} P.~Farah, A.~Popp, W.~A.~Wall: \emph{Segment-based vs.\ element-based integration for mortar methods in computational contact mechanics}, Computational Mechanics 55 (2015), 209--228. (Lokale Kopie: \path{meetings/literature/FarahPoppWall_2015_SegmentBasedVSElementBasedIntegrationForMortarMethodsInComputationalContactMechanics.pdf})

\bibitem[Gitterle et al.(2010)]{gitterle} M.~Gitterle, A.~Popp, M.~W.~Gee, W.~A.~Wall: \emph{Finite deformation frictional mortar contact using a semi-smooth Newton method with consistent linearization}, International Journal for Numerical Methods in Engineering 84 (2010), 543--571. (Lokale Kopie: \path{meetings/literature/GitterlePoppGeeWall_2010_FiniteDeformationFrictionalMortarContactUsingASemiSmoothNewtonMethodWithConsistentLinearization.pdf})

\bibitem[Hintermüller et al.(2002)]{hik2002} M.~Hintermüller, K.~Ito, K.~Kunisch: \emph{The primal-dual active set strategy as a semismooth Newton method}, SIAM Journal on Optimization 13 (2002), 865--888. (Lokale Kopie: \path{meetings/literature/HintermuellerItoKunisch_2002_ThePrimalDualActiveSetStrategyAsASemismoothNewtonMethod.pdf})

\bibitem[Hüeber \& Wohlmuth(2005)]{hueber2005} S.~H\"ueber, B.~I.~Wohlmuth: \emph{A primal-dual active set strategy for non-linear and contact problems}, Computer Methods in Applied Mechanics and Engineering 194 (2005), 3147--3166. (Lokale Kopie: \path{meetings/literature/HüeberWohlmuth_2005_APrimalDualActiveSetStrategyForNonLinearContactProblems.pdf})

\bibitem[Lamichhane et al.(2005)]{lamichhane2005} B.~P.~Lamichhane, R.~P.~Stevenson, B.~I.~Wohlmuth: \emph{Higher order mortar finite element methods in 3D with dual Lagrange multiplier bases}, Numerische Mathematik 102 (2005), 93--121. (Lokale Kopie: \path{meetings/literature/LamichhaneStevensonWohlmuth_2005_HigherOrderMortarFiniteElementMethodsIn3DWithDualLagrangeMultiplierBases.pdf})

\bibitem[Popp et al.(2009)]{poppgeewall2009} A.~Popp, M.~W.~Gee, W.~A.~Wall: \emph{A finite deformation mortar contact formulation using a primal--dual active set strategy}, International Journal for Numerical Methods in Engineering 79 (2009), 1354--1391. (Lokale Kopie: \path{meetings/literature/PoppGeeWall_2009_AFiniteDeformationMortarContactFormulationUsingAPrimalDualActiveSetStrategy.pdf})

\bibitem[Popp et al.(2010)]{popp2010} A.~Popp, M.~W.~Gee, W.~A.~Wall: \emph{A dual mortar approach for 3D finite deformation contact with consistent linearization}, International Journal for Numerical Methods in Engineering 83 (2010), 1428--1465.

\bibitem[Popp et al.(2012)]{popp2012} A.~Popp, B.~I.~Wohlmuth, M.~W.~Gee, W.~A.~Wall: \emph{Dual quadratic mortar finite element methods for 3D finite deformation contact}, SIAM Journal on Scientific Computing 34 (2012), B421--B446. (Lokale Kopie: \path{meetings/literature/PoppWohlmuthGeeWall_2012_DualQuadraticMortarFiniteElementMethodsFor3DFiniteDeformationContact.pdf})

\bibitem[Puso \& Laursen(2004)]{puso2004} M.~A.~Puso, T.~A.~Laursen: \emph{A mortar segment-to-segment contact method for large deformation solid mechanics}, Computer Methods in Applied Mechanics and Engineering 193 (2004), 601--629. (Lokale Kopie: \path{meetings/literature/PusoLaursen_2004_AMortarSegmentToSegmentContactMethodForLargeDeformationSolidMechanics.pdf})

\bibitem[Puso et al.(2008)]{puso2008} M.~A.~Puso, T.~A.~Laursen, J.~Solberg: \emph{A segment-to-segment mortar contact method for quadratic elements and large deformations}, Computer Methods in Applied Mechanics and Engineering 197 (2008), 555--566. (Lokale Kopie: \path{meetings/literature/PusoLaursenSolberg_2008_ASegmentToSegmentMortarContactMethodForQuadraticElementsAndLargeDeformations.pdf})

\bibitem[Sutherland \& Hodgman(1974)]{sutherland1974} I.~E.~Sutherland, G.~W.~Hodgman: \emph{Reentrant polygon clipping}, Communications of the ACM 17 (1974), 32--42. (Lokale Kopie: \path{meetings/literature/SutherlandHodgam_1974_ReentrantPolygonClipping.pdf})

\bibitem[Wohlmuth(2000)]{wohlmuth2000} B.~I.~Wohlmuth: \emph{A mortar finite element method using dual spaces for the Lagrange multiplier}, SIAM Journal on Numerical Analysis 38 (2000), 989--1012. (Lokale Kopie: \path{meetings/literature/Wohlmuth_2000_AMortarFiniteElementMethodUsingDualSpacesForTheLagrangeMutliplier.pdf})

\bibitem[Wriggers(2006)]{wriggers} P.~Wriggers: \emph{Computational Contact Mechanics}, 2.~Aufl., Springer, 2006. (Lokale Kopie: \path{meetings/literature/Wriggers_2006_ComputationalContactMechanics.pdf})

\end{thebibliography}

\end{document}
